% id: 123-e10
% book: 개념원리 공통수학1 · 15 이차방정식의 근과 계수의 관계
% section: 필수·발전 예제
% figure: none
% answer: ⑴ $3$ \quad ⑵ $a=2$, $b=1$
% answer_source: 본문 풀이
% variant_level: 0 (원본 전사)
% 이 파일은 build.mjs 가 items.json 에서 만든다. 고칠 때는 items.json 을 고친다.
\smash{\raisebox{1.5mm}{\dmkichul{개념원리 공통수학1 123쪽 예제 10 · 필수}}}
다음 물음에 답하시오.
\dmsubprob{1} 이차방정식 $x^2+ax+b=0$의 두 근이 $-4$, $2$일 때, 이차방정식 $ax^2+(a+b)x+b=0$의 두 근의 합을 구하시오. \cond{$a$, $b$는 상수이다.}
\dmsubprob{2} 이차방정식 $x^2-ax+b=0$의 두 근이 $\alpha$, $\beta$이고, 이차방정식 $x^2-(a+1)x+2=0$의 두 근이 $\alpha+\beta$, $\alpha\beta$일 때, 상수~$a$, $b$의 값을 구하시오.
